\chapter{مقدمه}

شکاف فزاینده بین سرعت پردازنده‌ها و حافظهٔ اصلی یکی از چالش‌های اساسی در معماری کامپیوتر است. برای پر کردن این شکاف از سلسله‌مراتب حافظهٔ نهان\footnote{\lr{Cache}} استفاده می‌شود. با این حال محدودیت‌هایی که در ساختارهای نگاشت حافظه وجود دارد، چالش‌های جدیدی را به وجود می‌آورد. به عنوان مثال، به دلیل سادگی و سرعت بالا در پیاده‌سازی، استفاده از حافظه‌های نهان با نگاشت مستقیم\footnote{\lr{Direct-Mapped}} یا انجمنی با درجهٔ پایین بسیار رایج است اما این روش باعث بروز نقصان‌های برخورد\footnote{\lr{Conflict Misses}} می‌شود. در این حالت ممکن است چندین آدرس پرکاربرد دقیقا به یک خط از حافظهٔ نهان اختصاص یابند. حتی اگر ظرفیت کلی حافظه برای نگهداری داده‌ها کافی باشد، این تداخل باعث می‌شود داده‌ها پیوسته از حافظه اخراج\footnote{\lr{Evict}} شوند.

این رفتار ناکارآمد که به عنوان پدیدهٔ تلاطم\footnote{\lr{Thrashing}} شناخته می‌شود، سامانه را مجبور می‌کند تا داده‌ها را بارها از سطوح پایین‌تر واکشی کند. این فرایند نه تنها ترافیک بیهوده‌ای در شبکهٔ ارتباطی حافظه ایجاد می‌کند بلکه تاخیر سنگینی را به پردازنده تحمیل می‌نماید. برای حل این مشکل، ایدهٔ اضافه کردن یک بافر بسیار کوچک و تمام‌انجمنی\footnote{\lr{Fully Associative}} به عنوان حافظهٔ نهان قربانی\footnote{\lr{Victim Cache}} مطرح می‌شود که داده‌های اخراجی از سطح اول را پیش از انتقال به سطح پایین‌تر در خود نگه می‌دارد. 

در این پروژه روند ناکارآمد ذکر شده ابتدا در قالب یک معماری پایه شبیه‌سازی و مدل‌سازی شده است تا میزان هدررفت عملکرد سامانه اندازه‌گیری شود. سپس با طراحی یک پیمانهٔ جدید و تغییر منطق مسیریابی پیام‌ها\footnote{\lr{Message Routing}}، سامانه بهینه شده است. هدف نهایی این گزارش تحلیل تاثیر این معماری جدید بر موازنهٔ میان هزینهٔ سخت‌افزاری و زمان دسترسی میانگین\footnote{\lr{Average Memory Access Time - AMAT}} برنامه‌ها و همچنین بررسی روند جلوگیری از بن‌بست\footnote{\lr{Deadlock}} در تبادل هم‌زمان پیام‌ها است.